% Commandes personnalisées LaTeX
% Fichier: config/commands.tex

% ============================================================================
% COMMANDES POUR LA GESTION DES IMAGES
% ============================================================================

% Commande principale pour l'inclusion d'images avec placeholder
\newcommand{\auditimage}[4][]{%
    % #1 = options (optionnel)
    % #2 = chemin de l'image
    % #3 = légende
    % #4 = label
    \IfFileExists{#2}{%
        % Image existante
        \begin{figure}[H]
            \centering
            \includegraphics[#1]{#2}
            \caption{#3}
            \label{fig:#4}
        \end{figure}
    }{%
        % Image manquante - placeholder
        \begin{figure}[H]
            \centering
            \placeholderbox{#2}{#3}{#4}
        \end{figure}
        \typeout{ATTENTION: Image manquante: #2}
    }
}

% Boîte de placeholder stylisée
\newcommand{\placeholderbox}[3]{%
    % #1 = chemin de l'image manquante
    % #2 = légende
    % #3 = label
    \fbox{%
    \begin{minipage}{0.8\textwidth}
        \centering
        \vspace{1cm}
        
        % Icône d'alerte
        {\Huge ⚠️}
        
        \vspace{0.5cm}
        
        % Informations sur l'image manquante
        \textbf{\textcolor{severitycritical}{IMAGE MANQUANTE}}\\
        \vspace{0.2cm}
        \texttt{#1}\\
        \vspace{0.5cm}
        
        % Détails de catégorie
        \small
        \begin{tabular}{p{0.3\textwidth} p{0.5\textwidth}}
            \textbf{Catégorie:} & \getcategory{#1} \\
            \textbf{Section:} & \getsection{#1} \\
            \textbf{Vulnérabilité:} & \getvulnerability{#1} \\
        \end{tabular}
        
        \vspace{0.5cm}
        
        % Légende
        \textit{#2}
        
        \vspace{1cm}
    \end{minipage}
    }
    \caption{#2 \textbf{(Placeholder - Image à ajouter)}}
    \label{fig:placeholder-#3}
}

% Commandes d'extraction d'informations depuis le chemin d'image
\newcommand{\getcategory}[1]{%
    \IfSubStr{#1}{/injection/}{Injection}{%
    \IfSubStr{#1}{/xss/}{XSS}{%
    \IfSubStr{#1}{/authentication/}{Authentification}{%
    \IfSubStr{#1}{/access_control/}{Contrôle d'accès}{%
    \IfSubStr{#1}{/configuration/}{Configuration}{%
    \IfSubStr{#1}{/data_exposure/}{Exposition de données}{%
    \IfSubStr{#1}{/methodology/}{Méthodologie}{%
    \IfSubStr{#1}{/other/}{Autre}{%
    Général}}}}}}}}
}

\newcommand{\getsection}[1]{%
    \IfSubStr{#1}{/sql/}{SQL Injection}{%
    \IfSubStr{#1}{/nosql/}{NoSQL Injection}{%
    \IfSubStr{#1}{/reflected/}{XSS Réfléchi}{%
    \IfSubStr{#1}{/dom/}{XSS DOM}{%
    \IfSubStr{#1}{/stored/}{XSS Stocké}{%
    \IfSubStr{#1}{/brute_force/}{Brute Force}{%
    \IfSubStr{#1}{/jwt/}{JWT}{%
    \IfSubStr{#1}{/idor/}{IDOR}{%
    \IfSubStr{#1}{/path_traversal/}{Path Traversal}{%
    \IfSubStr{#1}{/file_upload/}{File Upload}{%
    \IfSubStr{#1}{/csrf/}{CSRF}{%
    \IfSubStr{#1}{/ssrf/}{SSRF}{%
    \IfSubStr{#1}{/xxe/}{XXE}{%
    Général}}}}}}}}}}}}}
}

\newcommand{\getvulnerability}[1]{%
    \StrBefore{#1}{.}[\tempname]%
    \StrBehind{\tempname}{/}[\vulnname]%
    \vulnname%
}

% Commandes spécifiques par catégorie d'images
\newcommand{\injectionimage}[2][]{%
    \auditimage[#1]{images/injection/#2}{#2}{inj-#2}%
}

\newcommand{\xssimage}[3][]{%
    \auditimage[#1]{images/xss/#2/#3}{XSS #2: #3}{xss-#2-#3}%
}

\newcommand{\authimage}[3][]{%
    \auditimage[#1]{images/authentication/#2/#3}{Auth #2: #3}{auth-#2-#3}%
}

\newcommand{\accessimage}[3][]{%
    \auditimage[#1]{images/access_control/#2/#3}{Access #2: #3}{access-#2-#3}%
}

\newcommand{\configimage}[2][]{%
    \auditimage[#1]{images/configuration/#2}{Configuration: #2}{config-#2}%
}

\newcommand{\dataimage}[2][]{%
    \auditimage[#1]{images/data_exposure/#2}{Data Exposure: #2}{data-#2}%
}

\newcommand{\methodologyimage}[2][]{%
    \auditimage[#1]{images/methodology/#2}{Methodology: #2}{method-#2}%
}

% ============================================================================
% COMMANDES POUR LES VULNÉRABILITÉS
% ============================================================================

% Commande pour créer une entrée de vulnérabilité
\newcommand{\vulnerabilityentry}[5]{%
    % #1 = nom
    % #2 = catégorie OWASP
    % #3 = difficulté
    % #4 = sévérité
    % #5 = page
    #1 & \owaspcategory{#2} & \difficultystar{#3} & \severitylevel{#4} & #5 \\
}

% Commande pour les en-têtes de tableau de vulnérabilités
\newcommand{\vulnerabilityheader}{%
    \toprule
    \textbf{Vulnérabilité} & \textbf{Catégorie OWASP} & \textbf{Difficulté} & \textbf{Sévérité} & \textbf{Page} \\
    \midrule
}

% ============================================================================
% COMMANDES POUR LES RECOMMANDATIONS
% ============================================================================

% Commandes pour les priorités de recommandations
\newcommand{\prioritycritical}{\textcolor{severitycritical}{\textbf{CRITIQUE}}}
\newcommand{\priorityhigh}{\textcolor{severityhigh}{\textbf{ÉLEVÉE}}}
\newcommand{\prioritymedium}{\textcolor{severitymedium}{\textbf{MOYENNE}}}
\newcommand{\prioritylow}{\textcolor{severitylow}{\textbf{BASSE}}}

% Commande pour une recommandation priorisée
\newcommand{\priorityrecommendation}[2]{%
    \begin{tcolorbox}[
        colback=#1!5!white,
        colframe=#1!75!black,
        title=Recommandation (\textbf{#1}),
        fonttitle=\bfseries,
        breakable
    ]
    #2
    \end{tcolorbox}
}

% ============================================================================
% COMMANDES POUR LES PAYLOADS ET CODE
% ============================================================================

% Environnement pour les payloads dangereux
\newenvironment{dangerouspayload}{%
    \begin{tcolorbox}[
        colback=severitycritical!5!white,
        colframe=severitycritical!75!black,
        title=\textbf{Payload d'exploitation},
        fonttitle=\bfseries\small,
        breakable
    ]
    \begin{verbatim}
}{%
    \end{verbatim}
    \end{tcolorbox}
}

% Environnement pour les exemples de code sécurisé
\newenvironment{safecode}{%
    \begin{tcolorbox}[
        colback=successcolor!5!white,
        colframe=successcolor!75!black,
        title=\textbf{Code sécurisé},
        fonttitle=\bfseries\small,
        breakable
    ]
    \begin{verbatim}
}{%
    \end{verbatim}
    \end{tcolorbox}
}

% Commande pour afficher du code inline
\newcommand{\inlinecode}[1]{%
    \texttt{\textcolor{primarycolor}{#1}}%
}

% ============================================================================
% COMMANDES POUR LES RÉFÉRENCES ET CITATIONS
% ============================================================================

% Commande pour les références OWASP
\newcommand{\owaspref}[1]{%
    \textbf{OWASP #1}%
}

% Commande pour les références CVE
\newcommand{\cveref}[1]{%
    \textbf{CVE-#1}%
}

% Commande pour les références NIST
\newcommand{\nistref}[1]{%
    \textbf{NIST SP #1}%
}

% ============================================================================
% COMMANDES POUR LA MISE EN FORME DU TEXTE
% ============================================================================

% Commande pour le texte important
\newcommand{\important}[1]{%
    \textbf{\textcolor{primarycolor}{#1}}%
}

% Commande pour les termes techniques
\newcommand{\techterm}[1]{%
    \textit{\textcolor{secondarycolor}{#1}}%
}

% Commande pour les notes de bas de page de sécurité
\newcommand{\securityfootnote}[1]{%
    \footnote{\textcolor{dangercolor}{⚠️ #1}}%
}

% ============================================================================
% COMMANDES POUR LES STATISTIQUES ET MÉTRIQUES
% ============================================================================

% Commande pour afficher un pourcentage avec couleur
\newcommand{\percentwithcolor}[2]{%
    \textcolor{#1}{\textbf{#2\%}}%
}

% Commande pour les compteurs d'images
\newcommand{\imagecount}[1]{%
    \textbf{\getimagecount{#1}}%
}

\newcommand{\getimagecount}[1]{%
    % Cette commande serait implémentée dans un script Python
    \textbf{15} % Valeur temporaire
}

% Commande pour les pourcentages de complétude
\newcommand{\percentcomplete}[1]{%
    \percentwithcolor{successcolor}{85} % Valeur temporaire
}

% Commandes pour les totaux
\newcommand{\totalimagecount}{%
    \textbf{108} % Valeur temporaire
}

\newcommand{\totalpercentcomplete}{%
    \percentwithcolor{successcolor}{78} % Valeur temporaire
}

% ============================================================================
% COMMANDES POUR LES LISTES D'IMAGES MANQUANTES
% ============================================================================

% Commande pour lister les images manquantes
\newcommand{\listmissingimages}{%
    \begin{itemize}
        \item \texttt{images/injection/sql\_login.png}
        \item \texttt{images/xss/reflected/search.png}
        \item \texttt{images/authentication/brute\_force/hydra\_attack.png}
        % ... autres images manquantes
    \end{itemize}
    \practicetip{Exécutez le script \texttt{validate\_images.py} pour générer une liste complète des images manquantes.}
}

% ============================================================================
% COMMANDES POUR LA CONFIGURATION DES IMAGES
% ============================================================================

% Commande pour configurer tous les chemins d'images
\newcommand{\configureallimages}{%
    \graphicspath{%
        {images/}%
        {images/injection/}%
        {images/xss/}%
        {images/xss/reflected/}%
        {images/xss/dom/}%
        {images/xss/stored/}%
        {images/authentication/}%
        {images/authentication/brute\_force/}%
        {images/authentication/jwt/}%
        {images/authentication/password\_reset/}%
        {images/access\_control/}%
        {images/access\_control/idor/}%
        {images/access\_control/path\_traversal/}%
        {images/access\_control/file\_upload/}%
        {images/configuration/}%
        {images/data\_exposure/}%
        {images/other/}%
        {images/methodology/}%
        {images/methodology/tools/}%
        {images/methodology/environment/}%
        {images/assets/}%
    }%
}

% ============================================================================
% COMMANDES POUR LES RAPPORTS ET VALIDATIONS
% ============================================================================

% Commande pour générer un rapport de validation
\newcommand{\generatereport}{%
    \typeout{========================================}%
    \typeout{RAPPORT DE VALIDATION DU DOCUMENT}%
    \typeout{========================================}%
    \typeout{Document: Rapport d'Audit de Sécurité}%
    \typeout{Date: \today}%
    \typeout{Auteur: Auditeur de Sécurité}%
    \typeout{========================================}%
}

% ============================================================================
% COMMANDES UTILITAIRES
% ============================================================================

% Commande pour vider le cache des images
\newcommand{\cleargraphicscache}{%
    \pdfximage{}
}

% Commande pour forcer le chargement des images
\newcommand{\forceimageload}{%
    \immediate\write18{echo "Forcing image reload..."}
}

% ============================================================================
% COMMANDES DE DÉBOGAGE
% ============================================================================

% Commande pour afficher des informations de débogage
\newcommand{\debuginfo}[1]{%
    \typeout{[DEBUG] #1}%
}

% Commande pour vérifier l'existence d'un fichier
\newcommand{\checkfile}[1]{%
    \IfFileExists{#1}{%
        \typeout{[OK] Fichier trouvé: #1}%
    }{%
        \typeout{[ERREUR] Fichier manquant: #1}%
    }%
}

% ============================================================================
% INITIALISATION DES COMMANDES
% ============================================================================

% Configuration initiale des chemins d'images
\AtBeginDocument{%
    \configureallimages%
    \generatereport%
}

% Fin du fichier commands.tex